% ==== 图 2　双重指纹生成与账本更新流程图 ====
%
% 框里只写动作名，细节由说明书 [0040]–[0053] 承载——图是示意图，不是正文的
% 转抄。每个框的文字都短到一行放得下；需要两行的手写 \\，不要让它自动折行
% （自动折行会把中文断在词中间）。
%
% 主干用 positioning 串联（below=... of ...），插删一步只改相邻两个节点的
% below 关系。两条侧向通道：
%   S204 的「是」经右侧小框「取内容摘要值作规范化结果」后直接进 S206
%   x = 7.6  S203 与 S206 的「空」：汇流并入 S207，按「空」处理
% 分支说明用**方框**而不是浮动文字——浮动标注会压到汇流竖线上。
% ============================================
\begin{tikzpicture}[x=1cm,y=1cm,node distance=5mm]

  \node[pterm] (start) {一次已完成的工具调用};

  \node[pbox=8.0cm, below=of start] (s201)
    {S201\quad 生成调用指纹：工具名称＋规范化入参};

  \node[pbox=8.0cm, below=of s201] (s202)
    {S202\quad 预处理输出：剥离内部提示，去首尾空白};

  \node[pdia=2.4, below=of s202] (s203)
    {S203\\ 输出为空？};

  \node[pdia=2.4, below=of s203] (s204)
    {S204\\ 是引用式结果？};

  \node[pbox=8.0cm, below=of s204] (s205)
    {S205\quad 生成规范化结果\\
     失败：取错误签名\quad 成功：剔除易变字段};

  \node[pdia=2.4, below=of s205] (s206)
    {S206\\ 规范化结果为空？};

  \node[pbox=8.0cm, below=of s206] (gen) {生成证据指纹};

  \node[pbox=8.0cm, below=of gen] (s207)
    {S207\quad 更新账本\\
     新证据：连续计数清零\quad 重复或空：累加};

  \node[pbox=8.0cm, below=of s207] (s208) {S208\quad 更新同结果计数};

  \node[pbox=8.0cm, below=of s208] (s209) {S209\quad 输出元数据事件};

  \node[pterm, below=of s209] (end) {结束};

  % ---- 主干 ----
  \draw[pl] (start) -- (s201);
  \draw[pl] (s201)  -- (s202);
  \draw[pl] (s202)  -- (s203);
  \draw[pl] (s203)  -- (s204);
  \draw[pl] (s204)  -- (s205);
  \draw[pl] (s205)  -- (s206);
  \draw[pl] (s206)  -- (gen);
  \draw[pl] (gen)   -- (s207);
  \draw[pl] (s207)  -- (s208);
  \draw[pl] (s208)  -- (s209);
  \draw[pl] (s209)  -- (end);

  % ---- S204 的「是」：取内容摘要值，跳过 S205 直接进 S206 ----
  % 宽度必须放得下「取内容摘要值」这一整行（6 字约 2.3cm），窄了会把
  % 「作规范/化结果」断在词中间。
  \node[pbox=2.5cm, anchor=west] (take) at (4.5,0 |- s204)
    {取内容摘要值\\ 作规范化结果};
  \coordinate (m204) at ($(s205.south)!0.5!(s206.north)$);
  \draw[pl]  (s204.east) -- (take.west);
  \draw[pll] (take.south) -- (take.south |- m204) -- (m204);
  \node[plab, anchor=south west] at ($(s204.east)+(0.1,0.02)$) {是};

  % ---- S203 与 S206 的「空」：汇流后并入 S207 ----
  \draw[pl]  (s203.east) -- (8.1,0 |- s203) -- (8.1,0 |- s207) -- (s207.east);
  \draw[pll] (s206.east) -- (8.1,0 |- s206);
  \node[plab, anchor=south west] at ($(s203.east)+(0.1,0.02)$) {是};
  \node[plab, anchor=south west] at ($(s206.east)+(0.1,0.02)$) {是};
  \node[pnote, anchor=south east] at (8.0,0 |- gen) {按「空」处理};

  % ---- 「否」分支 ----
  \node[plab, anchor=west] at ($(s203.south)!0.5!(s204.north)+(0.12,0)$) {否};
  \node[plab, anchor=west] at ($(s204.south)!0.5!(s205.north)+(0.12,0)$) {否};
  \node[plab, anchor=west] at ($(s206.south)!0.5!(gen.north)+(0.12,0)$)  {否};

\end{tikzpicture}
